\section{The standard-form tensor and the actual local rule}
\label{sec:qca_tensor_local_rule}

The same shifted unitary must occur on both neighbouring output pairs.
The following identities put the four physical output coordinates of
\cite[equation~\textup{StandardForm}, lines~603--622]{Cirac2017MPU}
in the ordered pair coordinates of
\cite[equations~\textup{locMarg}, \textup{bruleU}, \textup{bruleV},
    lines~980--1008]{SchumacherWerner2004QCA}.

\begin{lemma}[Products in disjoint bipartite coordinates]
    \label{lem:qca_bipartite_kronecker_product}
    \lean{SpinChain.bipartiteMatrixToQuasiLocalObservable_kronecker}
    \notready
    \uses{lem:qca_bipartite_observable_embedding}
    Let $\Lambda$ and $\Gamma$ be disjoint finite regions. In the ordered
    configuration coordinates of $\Lambda\cup\Gamma$, the observable of
    $A\otimes B$ is the product of the observable of $A$ on $\Lambda$
    and the observable of $B$ on $\Gamma$.
\end{lemma}
\begin{proof}
    \notready
    Factor $A\otimes B=(A\otimes 1)(1\otimes B)$ and use the two
    canonical inclusions and preservation of multiplication.
\end{proof}

\begin{lemma}[Translation of an ordered physical pair]
    \label{lem:qca_consecutive_pair_observable_translation}
    \lean{SpinChain.quasiLocalTranslation_consecutivePair}
    \notready
    \uses{def:qca_consecutive_pair_coordinates}
    A matrix $A$ on the ordered pair $(a,a+1)$ is carried by translation
    through $t\in\mathbb Z$ to the same matrix on
    $(a+t,a+t+1)$. This includes negative $a$ and $t$.
\end{lemma}
\begin{proof}
    \notready
    Translate the two site coordinates and use the canonical inclusion
    along the equality of their finite regions.
\end{proof}

\begin{lemma}[Matrix-unit expansion of the four-site output]
    \label{lem:mpu_two_link_output_matrix_units}
    \lean{MPOTensor.twoLinkOutputObservable_eq_sum_kronecker}
    \notready
    \uses{def:mpu_two_site_open_three_block}
    Put $K=uAu^\dagger$, with its indices ordered as $(\ell,r)$.
    The four-site output matrix of the supplied standard form is
    \begin{align}
        \sum_{i,j}K_{ij}
        \bigl[v(1_r\otimes E_{i_1j_1})v^\dagger\bigr]
        \otimes
        \bigl[v(E_{i_2j_2}\otimes 1_\ell)v^\dagger\bigr].
    \end{align}
    Both shifted gates have intermediate order $(r,\ell)$. This identity
    uses no unitarity hypothesis.
\end{lemma}
\begin{proof}
    \notready
    Expand $K$ in matrix units inside $1_r\otimes K\otimes 1_\ell$,
    and distribute conjugation by $v\otimes v$ across the finite sum.
\end{proof}

\begin{lemma}[The tensor output equals the fixed local rule]
    \label{lem:qca_tensor_output_fixed_local_rule}
    \lean{SpinChain.PropagatesWithin.twoLinkOutputObservable_eq_constantMargolusGateRule}
    \notready
    \uses{lem:qca_bipartite_kronecker_product,
        lem:qca_consecutive_pair_observable_translation,
        lem:mpu_two_link_output_matrix_units,lem:qca_fixed_margolus_gates}
    Embed the four-site output in the ordered pairs $(-1,0)$ and $(1,2)$.
    Its observable is exactly the fixed Margolus expression for $K=uAu^\dagger$:
    the first shifted gate acts on $(-1,0)$ and the second on $(1,2)$.
\end{lemma}
\begin{proof}
    \notready
    Apply the preceding product and translation identities to every
    summand of the matrix-unit expansion.
\end{proof}

\begin{lemma}[A constructed tensor for the actual translated local rule]
    \label{lem:qca_constructed_tensor_actual_local_rule}
    \lean{SpinChain.PropagatesWithin.exists_standardFormTensor_actual_margolus_rule}
    \notready
    \uses{lem:qca_fixed_margolus_gates,def:mpu_tensor_from_two_gates,
        lem:mpu_tensor_from_gates_periodic_circuit,lem:qca_tensor_output_fixed_local_rule}
    Let $\omega$ be a translation-covariant automorphism with
    nearest-neighbour propagation on the homogeneous chain. There are
    unitary gates $u,v$ and a tensor $W$ on the ordered two-site alphabet,
    of positive bond dimension $dr$, with the supplied standard form for
    these gates. Its periodic operators are unitary at every positive
    length. For every $t\in\mathbb Z$ and every physical pair matrix $A$,
    \begin{align}
        \omega\bigl(\tau_t(A_{\{0,1\}})\bigr)
        =\tau_t\bigl(\iota_{\{-1,0\},\{1,2\}}
        (\mathcal O_{u,v}(A))\bigr),
    \end{align}
    where $\mathcal O_{u,v}(A)$ is the four-site matrix above and $\iota$
    uses the increasing physical coordinates of each pair. Identification
    with the completed action on the twice-blocked chain is a further step.
\end{lemma}
\begin{proof}
    \notready
    Use the fixed gates derived from the support algebras. Construct the
    tensor and its supplied standard form from those gates, then apply
    the preceding local-rule identity at every translate.
\end{proof}
